\section{Introduction}
\label{sec:introduction}

LoRa's long frames make mesh reliability sensitive to airtime,
radio configuration, and message length~\cite{augustin2016study,bor2016scale}.
A delivered message is not necessarily a confirmed transaction: a valid
ACK must also reach the source before the deadline. Forward delivery
and reverse confirmation therefore need separate measurements.

Reusing a known route confines forwarding. Missing confirmation, however,
does not identify whether DATA or ACK was lost, a relay was busy, or the
route failed. Repetition may overcome a transient loss; broader
forwarding may find another path. Either can instead add contention.
A recovery policy must bound additional work, admit it against remaining
time, and define when feedback can still cancel a queued transmission.

MeshEcho-CPR addresses these decisions with a finite controller. An
ACK-confirmed route is tried first when locally feasible; a routed flow
can then repeat once and attempt one recovery flood. A flow initially
served by a flood has neither later stage. A matching ACK can cancel a
relay copy only before that copy commits to the radio schedule. The source
also distinguishes a live message's confirmation from shared cache
evidence, because concurrent messages can confirm or invalidate routes
in a different order from their creation.

The research question is when this combination changes the trade-off
between timely source confirmation and physical transmit airtime, and
which component causes the change. Source routing, finite retries,
flooding, and ACK-assisted coordination have established antecedents
\cite{johnson1996dsr,liu2004truncated,biswas2005exor}. The contribution
here is a precise interaction of these mechanisms and an auditable
evaluation of its conditions of benefit and failure. A generic
retry-followed-by-flood sequence is not claimed as new.

We provide an executable-policy specification with
finite source-stage bounds, explicit state limits, and a defined
cancellation boundary. Our measurements retain every
scheduled flow and every physical transmission, plus clean controls for
the repeat, recovery, and ACK-cancellation components. A frozen 760-run
study with 20 paired seeds per condition exposes both the benefit of
recovery under path disruption and its substantial confirmation penalty
under higher traffic contention. The recovery control changes timely
confirmation in opposite directions: $+8.32$ percentage points in the
disrupted setting and $-15.99$ points in the contention setting. This
reversal, rather than a universal efficiency claim, is the principal
empirical finding.
The comparison concerns MeshEcho-CPR, legacy MeshEcho, Meshtastic-like,
and MeshCore-like behavioral models. The last two omit documented
product mechanisms; their results cannot establish performance superiority
over official firmware.
